\section{AC 自动机（Aho-Corasick）}

用于多模式串同时匹配：先逐个插入所有模式串，再 \texttt{build()} 建立自动机，之后在文本串上沿转移走一遍即可完成匹配，复杂度 $\mathcal{O}(\sum|\text{模式}|+|\text{文本}|)$。

\texttt{operator>>(s)} 插入一个模式串并返回其终止结点编号；\texttt{fail} 为失配指针、\texttt{in} 为 \texttt{fail} 树入度。精确统计每个模式串出现次数时，把文本串经过的结点计数沿 \texttt{fail} 树按 \texttt{in} 度做拓扑累加即可。

\begin{lstlisting}
template<int S,int E>
struct AhoCorasick{
    AhoCorasick(){
        nodes.resize(1);
    }
    static constexpr int charSiz = E-S+1;
    struct ACNode{
        int fa = 0, fail = 0, in = 0;
        bool built = 0;
        std::vector<int> nxt = std::vector<int>(charSiz);
        int& operator[](int i){return nxt[i];}
    };
    std::vector<ACNode> nodes;
    constexpr int ch(int raw){
        return raw - S;
    }
    constexpr int newNode(int fa = 0){
        nodes.push_back(ACNode());
        nodes.back().fa = fa;
        return nodes.size() - 1;
    }
    int operator>>(const std::string& s){
        assert(!nodes[0].built);
        int cur = 0;
        for (int i = 0;i < s.size();i ++){
            int chd = ch(s[i]);
            if (!nodes[cur][chd]) nodes[cur][chd] = newNode(cur);
            cur = nodes[cur][chd];
        }
        return cur;
    }
    int& fail(int u){
        return nodes[u].fail;
    }
    void build(){
        if (nodes.front().built) return;
        nodes.front().built = 1;
        std::queue<int> bfs;
        for (int v : nodes[0].nxt){
            if (v) {
                bfs.push(v);
                nodes[v].built = 1;
            }
        }
        while(bfs.size()){
            int u = bfs.front();
            bfs.pop();
            for (int i = 0;i < charSiz;i ++){
                const int &chd = nodes[u][i];
                if (chd) {
                    if (!nodes[chd].built){
                        nodes[chd].built = 1;
                        bfs.push(chd);
                        fail(chd) = nodes[fail(u)][i];
                        nodes[fail(chd)].in ++;
                    }
                } else {
                    nodes[u][i] = nodes[fail(u)][i];
                }
            }
        }
    }
    ACNode& operator[](int i){
        return nodes[i];
    }
};
using ACAM = AhoCorasick<'a','z'>;
\end{lstlisting}

\section{广义后缀自动机（GSAM）}

用于多个模式串共享的广义后缀自动机，支持多串的子串统计类问题（本质不同子串数、出现次数等）。\texttt{operator>>} 依次插入各模式串。

结点转移为 \texttt{ch}、后缀链接为 \texttt{link}、\texttt{longest} 为该状态最长串长、\texttt{cnt} 为该状态的结束次数；本质不同子串总数为 $\sum_i(\texttt{longest}[i]-\texttt{longest}[\texttt{link}[i]])$。转移表按固定字符集开定长数组，默认支持 'a'–'z'；若字符集不同或很大（约 50 以上），需自行把 \texttt{ch} 换成按需存储的稀疏表。

\begin{lstlisting}
struct SAMInfo{
    int cnt = 0;
    std::vector<int> ch;
    int link = 0,longest = 0;
    int& operator[](int idx){return ch[idx];}
};

// General SuffixAutomaton
template<typename Info,int S,int E>
struct SuffixAutomaton{
    std::vector<Info> nodes;
    int ed = 1;
    constexpr SuffixAutomaton(){
        newNode(); // Void Node
        newNode(); // Empty String Node -> ""
    }
    constexpr int newNode(){
        int idx = nodes.size();
        nodes.push_back(Info());
        nodes.back().ch.resize(E-S+1);
        return idx;
    }
    constexpr void extend(int c){
        if (nodes[ed][c] && nodes[ed].longest +1 == nodes[nodes[ed][c]].longest){
            ed = nodes[ed][c],nodes[ed].cnt ++;
            return;
        }
        int p = newNode(),lst = ed;
        nodes[p].longest = nodes[ed].longest + 1,nodes[p].cnt = 1;
        for (;ed && !nodes[ed][c];ed = nodes[ed].link) nodes[ed][c] = p;
        if (!ed) nodes[p].link = 1;
        else {
            int q = nodes[ed][c];
            if (nodes[q].longest == nodes[ed].longest + 1) nodes[p].link = q; // Legal link
            else {
                int x;
                if (lst == ed) x = p; // Split self
                else x = newNode(),nodes[p].link = x; // New Link
                nodes[x].ch = nodes[q].ch,nodes[x].link = nodes[q].link;
                nodes[x].longest = nodes[ed].longest + 1;
                nodes[q].link = x;
                for (;ed && nodes[ed][c] == q;ed = nodes[ed].link) nodes[ed][c] = x;
            }
        }
        ed = p; // update ed
    }
    constexpr void operator>>(const std::string &str){
        ed = 1;
        for (auto &c : str) extend(c-S);
    }
    constexpr Info& operator[](int idx){return nodes[idx];}
};
using GSAM = SuffixAutomaton<SAMInfo,'a','z'>;
\end{lstlisting}

\section{回文树（PAM / Eertree）}

用于求字符串的全部本质不同回文子串及其相关统计。除两根外每个结点对应一个本质不同回文，\texttt{len} 为其长度、\texttt{fail} 为其最长回文真后缀；\texttt{extend(c)} 在线加入一个字符并返回以新字符结尾的回文后缀个数（\texttt{cnt} 沿 \texttt{fail} 链累计）。

注意 \texttt{cnt} 不是出现次数，统计各回文出现次数需另行沿 \texttt{fail} 树拓扑累加。

\begin{lstlisting}
struct PAMNode{
    int fail = 1,len = 0,cnt = 0;
    std::vector<int> ch = std::vector<int>(26);
    int& operator[](int idx){return ch[idx];}
};
struct PalindromicAutomaton{
    using Info = PAMNode;
    std::vector<Info> nodes;
    int lst = 1;
    std::string raw;
    PalindromicAutomaton():nodes(2){
        nodes[0].fail = 1,nodes[1].len = -1;
    }
    int newNode(){
        nodes.push_back(Info());
        return nodes.size()-1;
    }
    int matchFail(int u,int idx){
        while(idx - nodes[u].len - 1 < 0 || raw[idx] != raw[idx-nodes[u].len-1]) u = nodes[u].fail;
        return u;
    }
    // 返回插入该字符后添加的回文子串数量
    int extend(int c){
        raw += c + 'a';
        int p = matchFail(lst,raw.size()-1);
        if (!nodes[p][c]){
            int x = newNode();
            nodes[x].fail = nodes[matchFail(nodes[p].fail,raw.size()-1)][c];
            nodes[p][c] = x;
            nodes[x].len = nodes[p].len + 2;
            nodes[x].cnt = nodes[nodes[x].fail].cnt + 1;
        }
        lst = nodes[p][c];
        return nodes[lst].cnt;
    }
};
using PAM = PalindromicAutomaton;
\end{lstlisting}
